\title{The Era of 1-bit LLMs: All Large Language Models are in 1.58 Bits}

\begin{abstract}
Recent advancements, including BitNet~\cite{bitnet}, are ushering in a new generation of 1-bit Large Language Models (LLMs). In this paper, we present a 1-bit LLM variant called \textbf{\text{BitNet b1.58}}, where each individual parameter (or weight) of the LLM is ternarized to \{-1, 0, 1\}. This model matches the performance of full-precision (i.e., FP16 or BF16) Transformer LLMs of the same scale and trained on the same token count, both in perplexity and downstream task results, while offering significant improvements in latency, memory usage, throughput, and energy efficiency. More importantly, the 1.58-bit LLM introduces a novel scaling law and training strategy for developing future generations of LLMs that are both efficient and high-performing. Additionally, it paves the way for a new computational framework and facilitates the design of specialized hardware tailored for 1-bit LLMs.
\end{abstract}

\section{The Era of 1-bit LLMs}

Over recent years, the AI community has witnessed rapid advancements in the size and capabilities of Large Language Models (LLMs). These models have achieved exceptional results across numerous natural language processing tasks; however, their growing scale presents deployment difficulties and raises concerns about environmental and economic impacts due to substantial energy demands.

One solution to these issues is post-training quantization, which generates low-bit models for inference~\cite{smoothquant, gptq, quip, quip_sharp}. This approach reduces the bit precision of weights and activations, thereby lowering the memory footprint and computational cost of LLMs. The evolution has progressed from 16-bit representations down to lower precision formats such as 4-bit variants~\cite{gptq, awq}. Nonetheless, despite its widespread industrial adoption, post-training quantization remains a suboptimal technique.

Recent research on 1-bit model architectures, such as BitNet~\cite{bitnet}, offers a promising approach to decrease the cost of LLMs while preserving their effectiveness. Standard LLMs operate with 16-bit floating-point values (e.g., FP16 or BF16), and the primary operation in these models is matrix multiplication. Consequently, the main computational expense arises from floating-point addition and multiplication operations. In contrast, BitNet’s matrix multiplication relies solely on integer addition, which drastically reduces the energy consumption for LLMs. Since power consumption is often the fundamental bottleneck in chip computation performance, this energy efficiency also enables faster processing.

Beyond computation, transferring model parameters from DRAM to the on-chip accelerator memory (such as SRAM) can be costly during inference. Expanding SRAM size to enhance throughput has been explored, but this leads to significantly higher costs compared to DRAM. Relative to full-precision models, 1-bit LLMs have a substantially smaller memory footprint both in terms of capacity and bandwidth. This reduction can greatly decrease the cost and duration of loading weights from DRAM, resulting in quicker and more efficient inference.

In this paper, we propose a notable 1-bit LLM variant named \textbf{\text{BitNet b1.58}}, where each parameter is ternary, taking values in \{-1, 0, 1\}. By adding the value 0 to the original 1-bit BitNet, we achieve an effective 1.58 bits in the binary system. \text{BitNet b1.58} preserves all the advantages of the original 1-bit BitNet, including its novel computational approach that almost eliminates multiplication in matrix multiplication and allows for high optimization. It also maintains the same energy consumption level as the original 1-bit BitNet and demonstrates significantly improved efficiency regarding memory use, throughput, and latency compared to FP16 LLM baselines. Moreover, \text{BitNet b1.58} provides two further benefits. First, its modeling capacity is enhanced due to explicit support for feature filtering enabled by the inclusion of 0 in the weights, which markedly boosts the performance of 1-bit LLMs. Second, our experimental results reveal that \text{BitNet b1.58} can achieve parity with full-precision (i.e., FP16) baselines in both perplexity and downstream task performance, starting at a 3B model size and using identical configurations (e.g., model size, training tokens, etc.).

\section{BitNet b1.58}

\text{BitNet b1.58} builds upon the BitNet framework, which is a Transformer model that substitutes \emph{nn.Linear} layers with \emph{BitLinear}. It is trained from the beginning, using weights quantized to 1.58 bits and activations quantized to 8 bits. Relative to the original BitNet, it incorporates several changes that we outline below.

\paragraph{Quantization Function.} To limit the weights to the set \{-1, 0, +1\}, we employ an \emph{absmean} quantization scheme. This method first normalizes the weight matrix by its mean absolute value, then rounds each element to the closest integer in \{-1, 0, +1\}:

\begin{equation}
    \widetilde{W} = \mathrm{RoundClip}\left(\frac{W}{\gamma + \epsilon}, -1, 1\right),
\end{equation}

\begin{equation}
    \mathrm{RoundClip}(x, a, b) = \max\left(a, \min\left(b, \mathrm{round}(x)\right)\right),
\end{equation}

\begin{equation}
    \gamma = \frac{1}{nm} \sum_{ij} |W_{ij}|.
\end{equation}

For activations, the quantization function follows the same approach as in BitNet, with the distinction that we do not rescale activations before applying nonlinearities into the $[0, Q_b]$ range. Instead, each token's activations are scaled to $[-Q_b, Q_b]$, thereby eliminating zero-point quantization. This simplifies both implementation and system-level optimization, while having a negligible impact on performance based on our experiments.

\paragraph{LLaMA-alike Components.}

The LLaMA architecture~\cite{llama, llama2} has become the standard foundation for open-source large language models. To align with the open-source ecosystem, our \text{BitNet b1.58} design incorporates components inspired by LLaMA. In particular, it employs RMSNorm~\cite{rmsnorm}, SwiGLU~\cite{swiglu}, rotary positional embeddings~\cite{rope}, and removes all bias terms. This design choice enables \text{BitNet b1.58} to be seamlessly integrated into widely-used open-source platforms (such as Huggingface, vLLM~\cite{vllm}, and llama.cpp\footnote{\url{https://github.com/ggerganov/llama.cpp}}) with minimal effort.

\section{Results}

We conducted a comparison between \text{BitNet b1.58} and our reproduced FP16 \text{LLaMA LLM} across multiple model sizes. To guarantee an equitable comparison, both models were pre-trained on the RedPajama dataset~\cite{redpajama} using 100 billion tokens. Their zero-shot performance was assessed on several language benchmarks, including ARC-Easy~\cite{arc}, ARC-Challenge~\cite{arc}, Hellaswag~\cite{hellaswag}, Winogrande~\cite{winoGrande}, PIQA~\cite{piqa}, OpenbookQA~\cite{openbookqa}, and BoolQ~\cite{boolq}. Additionally, we reported validation perplexity on the WikiText2~\cite{wikitext2} and C4~\cite{c4} datasets.

We also measured and compared the runtime GPU memory usage and latency of both \text{LLaMA LLM} and \text{BitNet b1.58}. These metrics were obtained using the FasterTransformer\footnote{\url{https://github.com/NVIDIA/FasterTransformer}} codebase, which is highly optimized for reducing LLM inference latency on GPU platforms. The 2-bit kernel from Ladder~\cite{ladder} was incorporated for \text{BitNet b1.58}. We focused on reporting the time per output token, as this represents the principal cost during inference.

Table~\ref{tab:ppl} presents a summary of perplexity and resource costs for \text{BitNet b1.58} and \text{LLaMA LLM}. The data indicate that \text{BitNet b1.58} achieves perplexity comparable to full precision \text{LLaMA LLM} starting at the 3B model size, while being 2.71 times faster and requiring 3.55 times less GPU memory. Notably, \text{BitNet b1.58} with a 3.9B model size runs 2.4 times faster, uses 3.32 times less memory, yet significantly outperforms \text{LLaMA LLM} 3B.

Table~\ref{tab:zero-shot} details the zero-shot accuracy on downstream tasks. We employed the evaluation pipeline from \emph{lm-evaluation-harness}\footnote{\url{https://github.com/EleutherAI/lm-evaluation-harness}} for this assessment. Results demonstrate that the performance gap between \text{BitNet b1.58} and \text{LLaMA LLM} diminishes as the model scale grows. Crucially, \text{BitNet b1.58} matches the full precision baseline’s performance starting at the 3B model size. Consistent with the perplexity findings, downstream task results reveal that \text{BitNet b1.58} 3.9B surpasses \text{LLaMA LLM} 3B while incurring lower memory and latency overheads. This confirms that \text{BitNet b1.58} constitutes a Pareto improvement over existing state-of-the-art LLM architectures.

\paragraph{Memory and Latency}

We additionally increased the model sizes to 7B, 13B, and 70B and assessed the associated costs. Figure~\ref{fig:latency-memory-energy} depicts the patterns of latency and memory usage, demonstrating that the acceleration improves as the model size expands. Notably, \text{BitNet b1.58} 70B achieves a speed-up that is 4.1 times greater than the \text{LLaMA LLM} baseline. This improvement arises because the time expense for \emph{nn.Linear} scales with the model size. Memory consumption exhibits a similar behavior, since the embedding remains in full precision and its relative memory share decreases for larger models. Both latency and memory metrics were obtained using a 2-bit kernel, indicating potential for further optimization to lower costs even more.

\paragraph{Energy}

We also evaluated the energy consumption of arithmetic operations for both \text{BitNet b1.58} and \text{LLaMA LLM}. Our primary focus was on matrix multiplication computations, as these dominate the cost in LLMs. Figure~\ref{fig:energy} shows the breakdown of energy costs. The majority of \text{BitNet b1.58}'s operations involve INT8 additions, whereas \text{LLaMA LLM} uses both FP16 additions and multiplications. Based on the energy model from~\cite{energycost, pokebnn}, \text{BitNet b1.58} achieves a 71.4-fold reduction in arithmetic operation energy consumption for matrix multiplication on 7nm chips. We also reported the end-to-end energy consumption for models processing 512 tokens. Our findings indicate that as the model size increases, \text{BitNet b1.58} becomes progressively more energy-efficient compared to the FP16 \text{LLaMA LLM} baseline. This efficiency gain is because the proportion of \emph{nn.Linear} computations rises with model size, while the energy cost from other components decreases for larger models.

\paragraph{Throughput}

We evaluate the throughput of \text{BitNet b1.58} and \text{LLaMA LLM} with 70B parameters using two 80GB A100 GPUs, employing pipeline parallelism~\cite{gpipe} to enable running the \text{LLaMA LLM} 70B on these devices. The batch size was increased until the GPU memory capacity was fully utilized, with a fixed sequence length of 512. As demonstrated in Table~\ref{tab:throughput}, \text{BitNet b1.58} 70B can handle up to 11 times the batch size of \text{LLaMA LLM}, leading to an 8.9-fold improvement in throughput.

\textbf{\text{BitNet b1.58} introduces a novel scaling law relating model performance and inference cost}. For reference, the following equivalences between model sizes in 1.58-bit and 16-bit formats are derived from the results shown in Figures~\ref{fig:latency-memory-energy} and \ref{fig:energy}.

\begin{itemize}
    \item A 13B BitNet b1.58 model is more efficient in terms of latency, memory consumption, and energy usage than a 3B FP16 LLM.
    \item A 30B BitNet b1.58 model outperforms a 7B FP16 LLM regarding latency, memory utilization, and energy efficiency.
    \item A 70B BitNet b1.58 model surpasses a 13B FP16 LLM in latency, memory demand, and energy consumption.
\end{itemize}

\paragraph{Training with 2T Tokens}

The quantity of training tokens is a critical parameter for LLMs. To assess the scalability of \text{BitNet b1.58} with respect to token count, we trained a \text{BitNet b1.58} model on 2 trillion tokens following the StableLM-3B~\cite{StableLM-3B-4E1T} data regimen, which is currently the leading open-source 3B model. Both models were tested on a benchmark comprising Winogrande~\cite{winoGrande}, PIQA~\cite{piqa}, SciQ~\cite{sciq}, LAMBADA~\cite{lambada}, and ARC-easy~\cite{arc}. Zero-shot accuracy results are presented in Table~\ref{tab:2t}. For tasks evaluated by accuracy and normalized accuracy, we averaged the two metrics. The StableLM 3B results at 2T tokens were sourced directly from its technical documentation. Our results indicate that \text{BitNet b1.58} delivers superior performance across all evaluated tasks, suggesting that 1.58-bit LLMs exhibit robust generalization capabilities.

\section{Discussion and Future Work}

\textbf{1-bit Mixture-of-Experts (MoE) LLMs}

Mixture-of-Experts (MoE) models have shown to be a cost-efficient strategy for large language models (LLMs). Although they substantially decrease computational FLOPs, their high memory demands and communication overhead between chips hinder broader deployment and practical use. These obstacles can be overcome by employing 1.58-bit LLMs. Firstly, the lowered memory requirements reduce the number of devices necessary for deploying MoE architectures. Additionally, this significantly diminishes the cost associated with transferring activations over networks. Ideally, if the entire model fits onto a single chip, communication overhead would be eliminated altogether.

\textbf{Native Support of Long Sequence in LLMs}

With the rise of LLMs, supporting long input sequences has become essential. A key challenge for long sequence inference lies in the memory usage caused by KV caches. \text{BitNet b1.58} marks a meaningful advancement toward native long sequence support by compressing activations from 16 bits down to 8 bits, effectively doubling the context length without increasing resource requirements. Further lossless compression down to 4 bits or below for 1.58-bit LLMs remains an avenue for future exploration.

\textbf{LLMs on Edge and Mobile}

Implementing 1.58-bit LLMs can substantially boost the performance of language models on edge and mobile platforms. These devices typically face constraints in memory capacity and computational power, limiting both the scale and effectiveness of LLMs. The reduced memory footprint and lower energy consumption of 1.58-bit LLMs make it feasible to deploy them on such devices, unlocking a variety of applications previously unattainable. This advancement can significantly enhance the functionalities of edge and mobile devices and foster innovative uses of LLM technology. Furthermore, 1.58-bit LLMs exhibit better compatibility with CPU architectures, which predominantly power edge and mobile hardware, enabling \text{BitNet b1.58} to run efficiently and further boost device capabilities.

\textbf{New Hardware for 1-bit LLMs}

Recent initiatives such as~Groq\footnote{\url{https://groq.com/}} have demonstrated encouraging outcomes and notable potential for developing dedicated hardware (e.g., LPUs) tailored for LLM workloads. Building on this progress, we propose and encourage efforts to create novel hardware and systems specifically optimized for 1-bit LLMs, leveraging the new computational framework introduced by BitNet~\cite{bitnet}. 

\end{document}